\DocumentMetadata{
  pdfversion=2.0,
  pdfstandard=ua-2,
  lang=en-US,
  tagging=on,
  tagging-setup={math/setup=mathml-SE,tabsorder=structure}
}
\documentclass[11pt,letterpaper]{article}
\usepackage[margin=0.68in,includefoot]{geometry}
\usepackage{newcomputermodern}
\usepackage{amsmath,amscd,mathtools}
\usepackage{tikz,adjustbox}
\providecommand{\square}{\mdlgwhtsquare}
\usepackage[hidelinks]{hyperref}
\hypersetup{
  pdftitle={Topology PhD Exam, May 2000},
  pdfauthor={University of Florida Department of Mathematics},
  pdfsubject={Graduate mathematics examination},
  pdfdisplaydoctitle=true,
  bookmarksopen=true
}
\setcounter{secnumdepth}{0}
\setlength{\parindent}{0pt}
\setlength{\parskip}{0.28em}
\setlength{\emergencystretch}{3em}
\setlength{\leftmargini}{2.15em}
\setlength{\leftmarginii}{2.65em}
\setlength{\leftmarginiii}{3em}
\renewcommand{\labelenumi}{\textbf{\theenumi.}}
\pagestyle{empty}
\raggedbottom
\allowdisplaybreaks
\newenvironment{examproblems}[1]{%
  \begin{enumerate}%
  \setcounter{enumi}{#1}%
  \setlength{\topsep}{0.35em}%
  \setlength{\partopsep}{0pt}%
  \setlength{\itemsep}{0.48em}%
  \setlength{\parsep}{0.18em}%
}{\end{enumerate}}
\newenvironment{examsubparts}{%
  \begin{enumerate}%
  \setlength{\topsep}{0.18em}%
  \setlength{\partopsep}{0pt}%
  \setlength{\itemsep}{0.16em}%
  \setlength{\parsep}{0.1em}%
}{\end{enumerate}}
\newenvironment{examinstructions}{%
  \par\begingroup\itshape
}{%
  \par\endgroup\vspace{0.18em}
}
\makeatletter
\renewcommand\section{\@startsection{section}{1}{0pt}%
  {-1.25ex plus -.25ex minus -.1ex}%
  {0.55ex plus .1ex}%
  {\normalfont\large\bfseries}}
\renewcommand{\@maketitle}{%
  \newpage\null\vskip -1em%
  {\footnotesize\bfseries
    \href{https://www.ufl.edu/}{University of Florida}, \href{https://math.ufl.edu/}{Department of Mathematics}\par}%
  \vskip -0.5em%
  \tagstructbegin{tag=Title}%
    {\LARGE\bfseries\href{https://gma.math.ufl.edu/exams/topology-phd/}{Topology PhD Exam}, May 2000\par}%
  \tagstructend%
  \vskip 0.25em%
  \tagmcbegin{artifact}%
    \hrule height 1pt%
  \tagmcend%
  \vskip 1em%
}
\makeatother
\title{Topology PhD Exam, May 2000}
\author{}
\date{}
\begin{document}
\maketitle
\thispagestyle{empty}
\begin{examinstructions}
Work the following problems and show all work. Support all statements to the best of your ability.
\end{examinstructions}
\begin{examproblems}{0}
\item
Show that there is no retraction of the Möbius band onto its boundary.
\item
State and prove the Contraction Mapping Theorem.
\item
Give the Mayer–Vietoris sequence for homology. Use this to compute the homology groups for the~\(n\)-sphere \(S^n\).
\item
Suppose that \(f,g:X\to S^n\) are two maps such that \(f(x)\) and \(g(x)\) are not antipodal for all \(x\in X\). Show that~\(f\) and~\(g\) are homotopic.
\item
Show that no covering space of the 2-torus has the homotopy type of a figure 8.
\end{examproblems}
\begin{examinstructions}
Answer the following with complete definitions or statements or short proofs.
\end{examinstructions}
\begin{examproblems}{5}
\item
Give a classification of the compact surfaces.
\item
State Urysohn’s Lemma. State the Tietze Extension Theorem.
\item
State the Five Lemma.
\item
Define the mapping cylinder of \(f:X\to Y\). Define the mapping torus of \(f:X\to X\).
\item
State the exact homology sequence for a pair of spaces.
\item
State the Jordan Curve Theorem.
\item
Define retraction and deformation retraction.
\item
State the Ascoli (Arzelà–Ascoli) Theorem.
\item
Let~\(V\) be a complete metric space such that no point is isolated. Can~\(V\) be a countably infinite set?
\item
Let~\(A\) be a countable subset of \(\mathbb{R}^2\). Show that \(\mathbb{R}^2\setminus A\) is arcwise connected.
\end{examproblems}
\end{document}
